\subsection{Реализация модели на языке Promela}

Модель системы находится в файле \code{modbus.pml}. В исходном файле
приведены комментарии к переменным и основным этапам обмена; полный код
доступен в репозитории проекта~\cite{repo}. Далее рассматриваются
функциональные части реализации. В листингах приведены фрагменты
исходного кода без комментариев, а из обработчиков также убрана часть
сообщений \code{printf}, предназначенных для просмотра обмена.

\subsubsection{Состояние системы и запуск процессов}

Для начала необходимо определить, какую информацию модель должна
сохранять между шагами. У Master есть состояние текущего обмена и поля
запроса. Память Slave представлена двумя массивами: индекс соответствует
адресу устройства. Нулевой элемент массивов не используется, поэтому
обращение к памяти Slave 2 записывается непосредственно через индекс 2.

\listing{model_data.pml}{Состояние Master и память устройств}

Все числовые поля объявлены как \code{byte}. При создании модели их
начальное значение равно нулю. После завершения транзакции содержимое
регистров сохраняется, поэтому следующие запросы работают с результатами
предыдущих записей. Флаг \code{slave_write_performed} показывает именно
факт выполнения операции: повторная запись уже хранящегося значения также
должна считаться записью.

Взаимодействие процессов организовано через каналы. Сообщение запроса
содержит адрес получателя, код функции, адрес регистра, значение и признак
CRC. В ответе последнее поле заменено кодом исключения. Нулевой код
исключения означает нормальный ответ. Процесс \code{init} создает три
экземпляра одного обработчика \code{Slave} и один экземпляр \code{Master}.

\listing{model_channels.pml}{Каналы обмена и начальный запуск}

Отдельные каналы запросов позволяют смоделировать получение общего кадра
всеми устройствами. Если бы все Slave читали один обычный канал Promela,
сообщение получил бы только один из них. В выбранной реализации адрес
проверяет сам обработчик, а не механизм доставки.

\subsubsection{Подготовка очередной транзакции}

Работа Master представляет собой бесконечное повторение одного обмена.
Алгоритм приведен на рисунке~\ref{fig:master}. Переход в
\code{MASTER_PREPARING} обозначает границу между транзакциями. Затем
Master очищает ответ и наблюдательные признаки предыдущего обмена,
сохраняет значения регистров и выбирает новый запрос.

\begin{figure}[htbp]
\centering
\includegraphics[width=0.76\linewidth]{master.png}
\caption{Алгоритм работы процесса Master}
\label{fig:master}
\end{figure}

Снимок памяти необходим для последующего сравнения. Проверка
«не было записи» и проверка «значение не изменилось» имеют разный смысл:
запись того же значения не меняет память, но операция при этом выполняется.
Поэтому используются и предыдущие значения, и отдельные флаги.

\listing{master_prepare.pml}{Сброс признаков обмена и сохранение памяти}

Далее Master выбирает параметры запроса. В конструкции \code{if} все
ветви с присваиваниями допустимы, поэтому выбор является
недетерминированным. Это задает множество вариантов поведения, которое
исследует SPIN, а не конкретную случайную последовательность запросов.

\listing{master_select.pml}{Выбор запроса и внесение одной ошибки}

Код \code{FC03} соответствует функции чтения регистров в Modbus,
но в рассматриваемой реализации эта функция специально не поддерживается.
Здесь она служит примером запроса, на который требуется вернуть исключение.
Проверка тем самым относится к устройству, реализующему только \code{FC06}.

После подготовки Master переходит в \code{MASTER_WAITING} и помещает
одинаковые поля в три канала. Сброс признаков, выбор параметров и отправка
объединены одним блоком \code{atomic}. Он сохраняет исключительное
выполнение последовательности, пока ее операторы не блокируются~\cite{atomic}.
К началу очередного обмена каналы запросов уже освобождены обработчиками.

\listing{master_send.pml}{Передача копий одного запроса}

\subsubsection{Проверки на стороне Slave и выполнение записи}

Каждый Slave работает в собственном бесконечном цикле и читает только
свой канал \code{slave_requests[slave_id - 1]}. Поля принятого кадра
сохраняются в локальных переменных. Поэтому один и тот же тип процесса
можно использовать для всех трех устройств, передавая ему собственный
адрес в параметре \code{slave_id}.

Алгоритм обработки показан на рисунке~\ref{fig:slave}. Сначала проверяется,
адресован ли запрос этому устройству или всем сразу. Затем проверяется CRC.
Только после этого определяется допустимость функции и адреса регистра.
Такой порядок позволяет игнорировать поврежденный кадр до выполнения
операции и формирования ответа.

\begin{figure}[htbp]
\centering
\includegraphics[width=0.98\linewidth,height=0.77\textheight,keepaspectratio]{slave.png}
\caption{Алгоритм обработки запроса одним Slave}
\label{fig:slave}
\end{figure}

В начале обработки \code{exception_code} равен нулю. При неподдерживаемой
функции он получает значение 1, при ошибочном адресе регистра --- 2.
Если обе проверки пройдены, Slave записывает выбранный регистр и
устанавливает \code{slave_write_performed[slave_id]}. Эти действия
объединены в \code{atomic}, поэтому состояние памяти и признак записи
наблюдаются согласованно.

\listing{slave_handler.pml}{Проверка запроса, запись и формирование ответа}

Вторая часть этого фрагмента определяет ответ. Для широковещательного
запроса выполняется \code{skip}: память может измениться, но сообщения
в канале ответов не возникает. При адресной успешной записи возвращаются
исходные поля запроса и нулевой код исключения.

Если обнаружена ошибка функции или регистра, к коду функции прибавляется
128, то есть устанавливается старший бит. В канале передаются также
нулевые значения полей регистра и данных и код исключения. Эти нули
относятся к фиксированной структуре сообщения модели: в ней нормальный
ответ и исключение используют один тип канала.

Отправка ответа и увеличение \code{responses_sent} также выполняются
совместно. Счетчик сохраняет факт отправки даже после того, как Master
извлек сообщение из канала. Поэтому свойство «не более одного ответа»
можно проверять на протяжении всей транзакции.

\subsubsection{Согласование завершения обмена}

После обработки запроса Slave устанавливает признак
\code{slave_processed[slave_id]}. Далее очищаются локальные поля и
отправляется служебное уведомление. Очистка локальных переменных не меняет память
устройства, но устраняет из состояния завершившего обработку процесса
данные старого кадра.

\listing{slave_done.pml}{Завершение обработки одного Slave}

Master читает три уведомления, после чего выбирает итоговое состояние.
При адресе 0 это \code{MASTER_BROADCAST_DONE}. В остальных случаях
проверяется наличие ответа: он извлекается в глобальные поля и
фиксируется \code{MASTER_RESPONSE}, либо выбирается \code{MASTER_TIMEOUT}.

\listing{master_complete.pml}{Получение результата транзакции}

Ожидание трех уведомлений существенно для проверки свойств. Оно позволяет
отличить отсутствие ответа от ситуации, когда нужный Slave еще не получил
возможность выполнить свой шаг. Кроме того, при переходе к следующей
транзакции все обработчики уже закончили предыдущую.

Для наблюдения за этими условиями дополнительно объявлены снимки памяти,
признаки завершения и счетчик ответов. Они используются в свойствах
системы, которые далее вынесены в отдельный файл.

\listing{model_observers.pml}{Данные для наблюдения за транзакцией}
\FloatBarrier
